\documentclass[11pt]{article}
\usepackage{amsmath,amssymb,hyperref,geometry}
\geometry{margin=1in}
\title{When a revised risk cap binds work already committed}
\author{}
\date{}
\begin{document}
\maketitle

\paragraph{Question and result.}
Can a warehouse dispatcher still admit an expedited load when the authority has changed the joint risk cap but the dispatcher has not received the change? For a cap on all unresolved physical exposure, moving the effective cut from issue to receipt does not preserve the fast option if its commitment and risk survive until receipt. A different rule that grandfathers the old commitment does preserve the option. An earlier receipt can recover service speed on other reached cuts, but cannot repair exposure from a reserve already locked to the wrong dispatch. These claims are conditional on the plant and service law below.

\paragraph{Inherited plant and fixed loss.}
Use the exact reserve model in \href{https://github.com/eternal-november/eos-research-programmes/pull/351}{PR 351}: two old loads are owed dispatches 3 and 4, one reserve protects exactly one dispatch, and its placement is irrevocable at the tick-2 lock. The two state coordinates flip independently at each step with probability $1/8$; the two admitted sensing laws vary report likelihood, not the physical transitions. The whole report word and the true tick-2 parity reach the dispatcher before lock. For each of the four report words, both parities and both laws have strictly positive mass. The remaining old loss $O$ is the physical loss event on the unprotected dispatch, as defined by the source loss cells. All copies of the experiment, including no-revision controls, add the same third duty: send a new load by tick 4. The inexpensive standard mode dispatches at 4 and adds no loss event. Fast mode dispatches at 3 and has a synthetic extra loss event $N$ of probability $1/20$, independent of $O$, reports, reserve selection and the authority branch conditional on the audit cut. The risk is $\Pr(O\cup N)$ through tick 4, not a sum of per-job caps; it is the same event definition and horizon under tightening, relaxation and no revision. Dispatch time is a service schedule, not a claim of damage-free arrival. The loss extension is a controlled model, not a fit to warehouse reliability data. No rerouting, moving the locked reserve, adding reserves, dropping an owed dispatch or changing the loss is an admissible mitigation.

\paragraph{Permission, binding cut, and observations.}
One authenticated safety authority can replace the scalar cap. Its scripted issue occurs after reserve lock and before the local admission phase. Receipt is either before admission, or after admission but before any physical dispatch. These phases are ordered, and neither mode observes a fresh plant state or loss outcome during them. The change/no-change branch is not seen before receipt and is independent of the hidden state given the cut; the controller has the report word, parity, own reserve placement, clock and received revisions, but neither the latent law label nor authority-side pending notices. Message receipt conveys the exact new cap and authority, and cannot be forged in this model. Admission at the stated phase fixes the mode and is the last chance for tick-3 dispatch. Waiting until a late receipt still permits standard dispatch by 4. Early-delivery controls fix a known receipt guarantee before admission for every issued message.

An issuer-effective joint rule binds at issue and checks already committed exposure then as well as all new exposure thereafter. Under a receipt-effective joint rule the old cap still governs admission before the receipt, but the new cap at receipt applies to every still-unresolved commitment, including the newly admitted load. Under the alternative receipt-effective \emph{admission-only} rule, the commitment already made under the old cap is grandfathered; it is not a promise that the revised cap holds for the old horizon. The policy cannot manufacture grandfathering for a stated joint rule. The scenario compares these as distinct contracts. It uses event cut plus authorized evidence, not a timestamp assigned retrospectively by the policy.

\paragraph{Exact residual exposure.}
The pre-issue lock is lawful under the old cap in both experiments. Tightening uses $25/64\to 1/5$; at the old cap an independent fair choice of either placement reaches both placements at the same report word and parity. Relaxation uses $1/5\to337/1024$, with the old parity policy; both parities and their distinct feasible exposures are reached at the old cap. Conditional on the report, parity and realized placement, the written source and independent trajectory enumeration give the following probabilities, identical in both laws:
\[
\begin{array}{c|c|c|c}
\text{parity}&\text{reserve}&\Pr(O)&\Pr(O\cup N)\text{ if fast}\\\hline
\text{equal}&3&175/1024&4349/20480\\
\text{equal}&4&400/1024&539/1280\\
\text{unequal}&3&337/1024&7427/20480\\
\text{unequal}&4&112/1024&197/1280
\end{array}
\]
Indeed $r+(1-r)/20$ is the fast risk for existing exposure $r$. More generally, for any law and audit cut with $r<1$, the exact requirement for admitting new loss is
\[
 \Pr(N\mid O^c,\text{audit cut})\le (c-r)/(1-r).
\]
A negative right-hand side means the existing physical exposure already fails: even standard mode cannot repair it. With tightening effective just before reserve lock, the parity policy instead places the reserve so standard mode can satisfy $1/5$ in every case. After lock, under tightening and identical equal-parity reports, reserve at 3 still permits standard service; reserve at 4 permits neither available mode. The same separation holds for unequal parity with the placements reversed; unequal-parity matched exposure leaves room for fast service. This failure does not require inferring the latent sensing law or any hidden mitigation assumption about dispatchers; it survives full information within this plant.

\paragraph{Transition theorem at the unnotified cut.}
For a fixed local observation $z$, let $H(z)$ contain all reached authority and physical histories consistent with that observation, including each admitted hidden law. For each $h$ let $A(h)$ be the actions satisfying the service deadline and every cap governing the commitments the action enters, including at a later binding cut when those commitments persist. A deterministic locally executable policy satisfies every audited history iff it chooses from $\bigcap_{h\in H(z)}A(h)$ at every reachable $z$: necessity is identical choice on identical observations, and sufficiency is membership in the actual history's feasible set. If audit is conditional on the realized action, a randomized policy must have support in the same intersection. This is not the criterion for a different contract auditing a policy mixture before the draw.

For equal parity and reserve at 3, old cap $25/64$ permits fast. A late tightening forbids fast under either joint rule: at issue for the issuer-effective rule, or when the outstanding fast commitment reaches the receipt audit for the receipt-effective joint rule. Before late receipt the dispatcher cannot distinguish the tightening and unchanged branches. Standard mode is the only robust choice, costing one tick on the unchanged branch compared with a dispatcher that knows the order status at admission. Under explicit receipt-effective grandfathering, fast is allowed at admission in both branches, but the pending joint exposure exceeds the tightened cap in the changed branch. This is not a workaround for either joint contract. For unequal parity and reserve at 4, fast remains legal after tightening; systematically waiting or choosing standard in both cuts wastes a tick there.

In the relaxation experiment, fast in the equal-parity cut is permitted after a timely verified relaxation, while no change keeps that cut standard-only. At the indistinguishable late admission, issuer-effectiveness alone cannot grant a fast policy safe across both branches. Receipt-effectiveness cannot grant it either: the old low cap still binds admission. The dispatcher can instead choose standard, or wait and lose the single fast opportunity. At the unequal-parity matched cut, fast is legal even before the relaxation, and waiting loses a tick. Revising the number does not move either reserve, and loosening a cap applicable only to new commitments cannot lower the actual inherited exposure. If auditing a previously high-cap committed, unmatched exposure at the intermediate cap, $337/1024$ would pass its old obligation alone and $400/1024$ would not; these are not lawful initial cuts under an already-effective $1/5$ rule.

\paragraph{Controls and limits of the comparison.}
At equal parity and reserve at 3, a new loss with marginal $1/20$ can be nested inside $O$, independent of $O$, or disjoint from $O$. The union risks are $175/1024$, $4349/20480$, and $1131/5120$ respectively. The first passes the tight cap and the other two do not. The marginal alone cannot decide admissibility. These are exact coupling controls, not claims that the nested or disjoint hazards are available interventions in the installed plant. Another control hides actual old placement and lets mismatches occur independently of the latent state with conditional probability $1/10$. Old risk pooled over placement is $395/2048$ or $269/2048$ depending on parity and passes $1/5$; it fails at each mismatched placement-conditioned cut. Likewise on the matched equal-parity cut, auditing the new action before a hidden independent fast-mode draw would allow a pooled fast probability through $596/849$ under tightening. That is not a solution when the completed admission record is an audit input.

The direct comparator uses the same exact risk table, effective-time checks and observation-class intersections. There is no claim of novelty in the finite intersection criterion, proof-work saving, contextual transfer, full v0.5 acceptance, or operational certification. The checked benefit of timely delivery is one tick on the named cuts, not a latency cost estimate for other workloads. If dispatch happens before receipt, new state evidence arrives before the effective cut, the loss horizon changes, or the legally binding event is generation rather than verified receipt, the transition conditions must be re-derived rather than reusing the result as stated.
\end{document}
